\chapter{Drift Catalogue and State Coverage}
\label{chap:drift_catalogue}

This page summarizes the synthetic drift space used in the evaluation. The
catalogue separates the temporal shape of a drift from its subtype. The mapping
to state types is intentionally approximate: a drift can leave secondary traces
in other state views, but the visual map marks the feature family where the
signal is expected to be most direct.

\begin{table}[ht]
\caption{Generic drift types and optional parameters.}
\label{tab:drift_types}
\centering
\small
\setlength{\tabcolsep}{3pt}
\begin{tabular}{@{}p{0.20\textwidth}p{0.46\textwidth}p{0.23\textwidth}@{}}
\toprule
\textbf{Drift type} & \textbf{Meaning} & \textbf{Parameters} \\
\midrule
\texttt{sudden} & Hard switch at the change point. & \texttt{change\_point} \\
\texttt{gradual\_linear} & Old and new behavior overlap; probability of the new behavior increases linearly. & \texttt{change\_point} \\
\texttt{gradual\_exponential} & Old and new behavior overlap; the new behavior rises slowly at first and faster near the end. & \texttt{change\_point} \\
\texttt{incremental} & Several intermediate versions are visited during the transition. & \texttt{change\_point}, \texttt{num\_versions} \\
\texttt{recurring} & Old and new behavior alternate for a finite period before the new behavior remains active. & \texttt{change\_point}, \texttt{period\_days} \\
\bottomrule
\end{tabular}
\end{table}

\begin{table}[ht]
\caption{Drift subtypes and subtype-specific parameters.}
\label{tab:drift_subtypes}
\centering
\scriptsize
\setlength{\tabcolsep}{2pt}
\begin{tabular}{@{}p{0.16\textwidth}p{0.22\textwidth}p{0.37\textwidth}p{0.15\textwidth}@{}}
\toprule
\textbf{Perspective} & \textbf{Subtype} & \textbf{Generated change} & \textbf{Parameters} \\
\midrule
Control-flow & \texttt{tree\_mutation} & Activities are added, deleted, moved, or process-tree operators are changed. & \texttt{change\_proportion}, \texttt{intensity} \\
\midrule
Resource & \texttt{reassignment} & Selected activities receive new dominant resources while preserving the 80/20 assignment shape. & \texttt{activities} \\
Resource & \texttt{workload\_distribution} & Event counts are shifted toward selected or sampled heavy resources. & \texttt{resources} \\
Resource & \texttt{service\_time} & Selected resources become faster or slower through duration multipliers. & \texttt{resources} \\
Resource & \texttt{pool\_size} & The active resource pool expands or contracts. & -- \\
Resource & \texttt{handover} & A dominant resource-to-resource handover target is rewired. & -- \\
\midrule
Inter-case & \texttt{arrival\_rate} & The mean time between case arrivals changes. & -- \\
Inter-case & \texttt{burstiness} & Arrival times become bursty or return to exponential spacing. & -- \\
Inter-case & \texttt{case\_mix} & The distribution of case types changes. & -- \\
Inter-case & \texttt{concurrency} & Marks a concurrency change; concurrency is observed indirectly through overlap. & -- \\
\midrule
Data & \texttt{numeric} & The mean and variance of \texttt{case:amount} change. & -- \\
Data & \texttt{categorical} & The distribution of \texttt{case:region} changes. & -- \\
\bottomrule
\end{tabular}
\end{table}

\begin{figure}[ht]
\centering
\begin{tikzpicture}[
    >={Latex[length=2mm]},
    every node/.style={font=\sffamily\small},
    box/.style={draw, rounded corners=2pt, line width=0.4pt, align=center,
        inner sep=5pt, minimum height=0.85cm, text width=2.8cm},
    source/.style={box, fill=gray!8},
    state/.style={box, fill=gray!18, text width=3.1cm},
    arrow/.style={->, line width=0.5pt},
    weak/.style={->, dashed, line width=0.4pt}
]
\node[source] (cf) {Control-flow\\\texttt{tree\_mutation}};
\node[source, below=0.55cm of cf] (data) {Data\\\texttt{numeric}, \texttt{categorical}};
\node[source, below=0.55cm of data] (res) {Resource\\assignment, workload, time};
\node[source, below=0.55cm of res] (inter) {Inter-case\\arrival, mix, overlap};

\node[state, right=2.2cm of data] (intra) {Intra-case state\\activity windows, gaps, case attributes};
\node[state, right=2.2cm of res] (resource) {Resource state\\workload, duration, handover};
\node[state, right=2.2cm of inter] (interstate) {Inter-case state\\arrivals, completions, concurrency};

\draw[arrow] (cf.east) -- (intra.west);
\draw[arrow] (data.east) -- (intra.west);
\draw[arrow] (res.east) -- (resource.west);
\draw[arrow] (inter.east) -- (interstate.west);
\draw[weak] (res.south east) to[bend left=8] (interstate.west);
\draw[weak] (inter.north east) to[bend right=8] (resource.east);
\end{tikzpicture}
\caption{Approximate relation between drift perspectives and the three state
types. Solid arrows show the primary state view; dashed arrows indicate common
secondary effects.}
\label{fig:drift_state_mapping}
\end{figure}

\begin{figure}[p]
\centering
\begin{tikzpicture}[
    every node/.style={font=\sffamily\scriptsize},
    statehead/.style={draw, rounded corners=3pt, line width=0.45pt,
        align=center, inner sep=5pt, minimum height=0.75cm,
        text width=0.27\textwidth, font=\sffamily\small\bfseries},
    feature/.style={draw, rounded corners=2pt, line width=0.35pt,
        align=left, inner sep=4pt, text width=0.27\textwidth},
    intrahead/.style={statehead, fill=blue!12, draw=blue!45!black},
    reshead/.style={statehead, fill=green!12, draw=green!45!black},
    interhead/.style={statehead, fill=orange!14, draw=orange!55!black},
    intrafeat/.style={feature, fill=blue!4, draw=blue!35!black},
    resfeat/.style={feature, fill=green!4, draw=green!35!black},
    interfeat/.style={feature, fill=orange!5, draw=orange!45!black},
    secondary/.style={feature, dashed, fill=gray!3, draw=gray!60},
    legend/.style={draw, rounded corners=2pt, fill=gray!5, align=left,
        inner sep=4pt, text width=0.83\textwidth}
]
\node[intrahead] (ih) {Intra-case state};
\node[reshead, right=0.48cm of ih] (rh) {Resource state};
\node[interhead, right=0.48cm of rh] (eh) {Inter-case state};

\node[intrafeat, below=0.28cm of ih] (i1) {
    \textbf{Activity-window vectors / local order}\\
    \texttt{tree\_mutation}
};
\node[intrafeat, below=0.20cm of i1] (i2) {
    \textbf{Case numeric attributes}\\
    \texttt{numeric}
};
\node[intrafeat, below=0.20cm of i2] (i3) {
    \textbf{Case categorical attributes}\\
    \texttt{categorical}; \texttt{case\_mix} if case type is used
};
\node[secondary, below=0.20cm of i3] (i4) {
    \textbf{Secondary temporal signal}\\
    \texttt{service\_time}: trace-local gaps can change
};

\node[resfeat, below=0.28cm of rh] (r1) {
    \textbf{Events per resource}\\
    \texttt{workload\_distribution}; \texttt{pool\_size}; \texttt{reassignment}
};
\node[resfeat, below=0.20cm of r1] (r2) {
    \textbf{Activity-resource event counts}\\
    \texttt{reassignment}; \texttt{workload\_distribution}
};
\node[resfeat, below=0.20cm of r2] (r3) {
    \textbf{Mean duration and wait per resource}\\
    \texttt{service\_time}
};
\node[resfeat, below=0.20cm of r3] (r4) {
    \textbf{Active resource/case load}\\
    \texttt{pool\_size}; \texttt{workload\_distribution}
};
\node[resfeat, below=0.20cm of r4] (r5) {
    \textbf{Resource handover counts}\\
    \texttt{handover}
};
\node[secondary, below=0.20cm of r5] (r6) {
    \textbf{Secondary load signal}\\
    \texttt{concurrency}: resource load can react indirectly
};

\node[interfeat, below=0.28cm of eh] (e1) {
    \textbf{New arrivals / total events}\\
    \texttt{arrival\_rate}
};
\node[interfeat, below=0.20cm of e1] (e2) {
    \textbf{Inter-event gap mean and spread}\\
    \texttt{burstiness}
};
\node[interfeat, below=0.20cm of e2] (e3) {
    \textbf{Case-type distribution}\\
    \texttt{case\_mix}
};
\node[interfeat, below=0.20cm of e3] (e4) {
    \textbf{Active, stalled, and completed cases}\\
    \texttt{concurrency}
};
\node[secondary, below=0.20cm of e4] (e5) {
    \textbf{Secondary throughput signal}\\
    \texttt{service\_time}; \texttt{tree\_mutation}
};

\node[legend, below=0.55cm of r6] {
    Filled cards indicate the most direct feature family for a subtype. Dashed
    cards mark secondary effects that may appear depending on the generated log
    and the chosen window size.
};
\end{tikzpicture}
\caption{Visual feature map from drift subtypes to the feature families used by
the three state views.}
\label{fig:drift_feature_mapping}
\end{figure}
